\section{Introduction}
\label{sec:intro}
Novel View Synthesis (NVS) is a fundamental task in 3D vision, enabling applications in virtual reality (VR), augmented reality (AR), and cinematic production. Neural implicit representations~\cite{mildenhall2020nerf, barron2021mip, barron2022mip, fridovich2022plenoxels, yu2021plenoctrees, muller2022instant, cao2023hexplane} have shown impressive progress, yet their training and rendering remain computationally expensive. Recently, 3D Gaussian Splatting (3DGS)~\cite{kerbl20233d} has emerged as a state-of-the-art method to achieve both improved rendering quality and real-time rendering speed. Although 3DGS enables real-time rendering and fast reconstructions, it still requires substantial training time, often several minutes or even hours per scene for reconstruction, since it relies on per-scene optimization. To address this limitation, several feed-forward-based approaches~\cite{charatan2024pixelsplat, szymanowicz2024splatter, chen2024mvsplat, xu2025depthsplat, kang2025selfsplat, ye2024no, huang2025no, zhang2024gs} have been proposed to directly predict the properties of 3D Gaussians (primitives) from multi-view images, leveraging large-scale datasets. In general, these approaches unproject the 2D pixels of each input image into a corresponding 3D primitive (i.e., \textit{pixel-aligned} Gaussians) and aggregate the resulting primitives from all input views to represent the corresponding scene.

After the feed-forward reconstruction stage, the predicted 3D Gaussian primitives can be stored and reused for real-time rendering on end-user devices such as mobile phones or AR/VR headsets. In practical applications, reconstruction is typically performed once on the server-side devices, while the generated 3DGS must support multiple rendering requests across diverse end-user devices. However, existing feed-forward methods typically generate a large number of pixel-aligned primitives, especially for dense-view reconstructions, which can be impractical for resource-constrained devices. To mitigate this, several approaches~\cite{jiang2025anysplat, zhang2024gaussian, ziwen2025long} have attempted to reduce redundant Gaussians across multiple views through voxelization~\cite{jiang2025anysplat}, graph-based aggregation~\cite{zhang2024gaussian}, or opacity-based pruning~\cite{ziwen2025long}. Nonetheless, these methods \textit{lack explicit control} over the target numbers of primitives \textit{because they were not specifically designed for fine-level control over this number}, which is a crucial factor for flexible deployment in real-time and streamable applications. Their threshold-based pruning strategies (e.g., voxel size or opacity-based pruning) often yield scene-dependent and unstable primitive counts, resulting in a poor quality-efficiency trade-off. In practice, \textit{explicit} control over the number of 3D Gaussian primitives is essential: without a mechanism to enforce a prescribed count, the final number of primitives varies across scenes and view counts, leading to unpredictable latency, memory, and bandwidth requirements. \textit{Importantly}, it should be noted that our work is distinguished from the standardization activities of MPEG~\cite{ISO2025} and other academic works~\cite{song2025tinysplat, chen2024fast} where 3D Gaussians are subject to \textit{compression} via codecs while our approach is related to effective and controllable \textit{compaction} of feed-forward 3DGS representations in terms of the number of the primitives.

To overcome these challenges, we propose EcoSplat, an \underline{E}fficiency-\underline{co}ntrollable feed-forward 3D Gaussian \underline{Splat}ting framework for dense-view reconstruction. As shown in Fig.~\ref{fig:figure_page1}, EcoSplat maximizes rendering quality while explicitly controlling the number of 3D Gaussians to match a target number, all in a feed-forward manner. To achieve this, we introduce a two-stage training process consisting of Pixel-aligned Gaussian Training (PGT) and Importance-aware Gaussian Finetuning (IGF). In the first stage (PGT), we train our model to reliably predict pixel-aligned 3D Gaussians. In the second stage (IGF), we further train our model to adaptively predict the 3D Gaussian parameters conditioned on the target number of primitives. Specifically, in the IGF stage, we introduce an importance-aware opacity loss $\mathcal{L}_\text{io}$ to encourage the suppression of less informative primitives. Moreover, for improved training stability, we adopt a Progressive Learning on Gaussian Compaction (PLGC) strategy that gradually expands the sampling range of the target number of primitives during training. As a result of this training process, our EcoSplat can adaptively select the most important Gaussians for any target number of primitives during inference (Fig.~\ref{fig:figure_page1}-(b)). Our main contributions are summarized as follows:
\begin{itemize}
  \setlength\itemsep{0.1cm}
  \item We introduce EcoSplat, the \textit{first} feed-forward 3DGS framework that reconstructs a high-quality and robust set of Gaussian primitives while satisfying \textit{explicit and controllable constraints on the target number of primitives}.
  \item We design a two-stage training process, including {Pixel-aligned Gaussian Training (PGT)} followed by {Importance-aware Gaussian Finetuning (IGF)}, that enables the model to learn an \textit{importance-aware primitive ranking} and achieve \textit{controllable efficiency}.
  \item We demonstrate state-of-the-art efficiency-controllable NVS performance on dense-view benchmarks such as RealEstate10K~\cite{zhou2018stereo} and ACID~\cite{infinite_nature_2020}. By achieving an optimal trade-off between primitive count and rendering quality, EcoSplat significantly outperforms previous feed-forward methods that lack fine-grained control over the number of primitives.
\end{itemize}



% As the number of input views increases, a large number of redundant Gaussians tend to be generated due to overlapping scene regions, leading to inefficiency in rendering~\cite{zhang2024gaussian}. 
% To tackle this limitation, several approaches~\cite{jiang2025anysplat, zhang2024gaussian, ziwen2025long, wang2025zpressor} have been proposed to reduce redundant Gaussians across multiple views, including voxelization~\cite{jiang2025anysplat}, graph-based aggregation~\cite{zhang2024gaussian}, and opacity-based pruning~\cite{ziwen2025long}. However, none of these methods provides explicit control over the final number of Gaussians. 